\documentclass[11pt]{article}

\usepackage[margin=1in]{geometry}
\usepackage{amsmath,amssymb,amsthm}
\usepackage{array}
\usepackage{booktabs}
\usepackage{enumitem}
\usepackage{hyperref}

% Common notation
\newcommand{\R}{\mathbb{R}}
\newcommand{\C}{\mathbb{C}}
\newcommand{\D}{\mathcal{D}}
\newcommand{\Sch}{\mathcal{S}}
\newcommand{\E}{\mathcal{E}}
\newcommand{\Lloc}{L^1_{\mathrm{loc}}}
\newcommand{\Cc}{C_c^\infty}
\newcommand{\pair}[2]{\langle #1, #2 \rangle}

\title{APM5557: Partial Differential Equations\\Functions, Equivalence Classes, and Distributions}
\author{}
\date{}

\begin{document}

\maketitle

The three kinds of object are related by a chain of maps:
\[
  \text{ordinary functions} \;\longrightarrow\; \text{equivalence classes}
  \;\longrightarrow\; \text{distributions}.
\]
\textbf{Each map forgets some information and makes more operations possible.
In particular, every object in the chain becomes differentiable once you reach
the distributions.}

\section{Ordinary functions}

\begin{description}[leftmargin=0pt]
  \item[What they are.] A rule $f:\Omega\to\R$ (or $\C$) with a definite value
    $f(x)$ at every point.
  \item[Equality.] $f=g$ means $f(x)=g(x)$ for \emph{every} $x$.
  \item[Typical spaces.] \leavevmode
    \begin{itemize}
      \item $C(\overline\Omega)$, with $\|f\|_\infty=\sup|f|$
      \item $C^k(\overline\Omega)$, with
        $\|f\| = \sum_{|\alpha|\le k}\sup|\partial^\alpha f|$
      \item $C^{k,\alpha}$, the H\"older spaces, and Lipschitz functions
      \item $C^\infty(\Omega)$, and $\Cc(\Omega)$, the smooth functions with
        compact support (the test functions)
    \end{itemize}
  \item[Language.] ``$f$ is continuous at $x_0$,'' ``$f(x_0)=3$,'' ``$f'(x_0)$
    exists.'' Evaluating at a point, $f\mapsto f(x_0)$, is a continuous linear
    functional on these spaces.
  \item[Derivatives.] Classical, defined pointwise as limits of difference
    quotients, and they may fail to exist.
\end{description}

\section{Equivalence classes (functions defined up to measure zero)}

\begin{description}[leftmargin=0pt]
  \item[What they are.] $[f]=\{g : g=f \text{ a.e.}\}$. Usually you just write
    $f$ and call it a \emph{representative}.
  \item[Why we need them.] $\int|f|^p=0$ implies only $f=0$ a.e., not
    $f\equiv 0$. On functions, $\|\cdot\|_p$ is therefore just a seminorm. It
    becomes a true norm only after you identify functions that agree a.e.
    Formally,
    \[
      L^p = \mathcal{L}^p \big/ \{f : f=0 \text{ a.e.}\}.
    \]
  \item[Equality.] $f=g$ in $L^p$ means $f=g$ almost everywhere.
  \item[Typical spaces.] \leavevmode
    \begin{itemize}
      \item $L^p(\Omega)$ for $1\le p\le\infty$; $L^\infty$ uses the essential
        supremum
      \item $\Lloc(\Omega)$, the functions integrable on every compact subset
      \item The Sobolev spaces $W^{k,p}(\Omega)$ and $H^k=W^{k,2}$, with
        $H^1_0\subset H^1$
    \end{itemize}
  \item[Language.] \leavevmode
    \begin{itemize}
      \item ``$f(x_0)$'' is \textbf{meaningless}, because a single point has
        measure zero and you can change $f$ there freely.
      \item ``$u\in H^1(a,b)$ is continuous'' means that $u$ \textbf{has a
        continuous representative}, and you agree to use it. Since two
        continuous functions that agree a.e.\ agree everywhere, that
        representative is unique, which is why we can treat $C$ as a subspace
        of $L^p$.
      \item Boundary values such as ``$u=0$ on $\partial\Omega$'' in higher
        dimensions need \textbf{trace theory}, because $\partial\Omega$ has
        measure zero.
    \end{itemize}
  \item[Derivatives.] Weak derivatives. $g=u'$ means
    \[
      \int u\,\varphi' = -\int g\,\varphi \quad\text{for every } \varphi\in\Cc.
    \]
    The result is again an equivalence class.
\end{description}

\section{Distributions}

\begin{description}[leftmargin=0pt]
  \item[What they are.] Continuous linear functionals $T:\Cc(\Omega)\to\R$.
    They aren't functions of $x$ at all. They are defined only through how they
    act on test functions, written $\pair{T}{\varphi}$.
  \item[Equality.] $T=S$ means $\pair{T}{\varphi}=\pair{S}{\varphi}$ for every
    test function $\varphi$.
  \item[Typical spaces.] \leavevmode
    \begin{itemize}
      \item $\D'(\Omega)$, the dual of $\Cc(\Omega)$
      \item $\Sch'(\R^n)$, the tempered distributions (the dual of the Schwartz
        space), which is where the Fourier transform lives
      \item $\E'(\Omega)$, the distributions with compact support
      \item Negative-order Sobolev spaces such as $H^{-1}=(H^1_0)'$, and $H^s$
        for $s<0$
    \end{itemize}
    Spaces like $\D'$ and $\Sch'$ are \textbf{not normed}. They carry weak-*
    topologies: $T_n\to T$ means $\pair{T_n}{\varphi}\to\pair{T}{\varphi}$ for
    every $\varphi$. Spaces like $H^{-1}$, by contrast, are Hilbert spaces.
  \item[How functions fit in.] Each $f\in\Lloc$ defines a distribution by
    \[
      \pair{T_f}{\varphi} = \int f\,\varphi .
    \]
    By the fundamental lemma of the calculus of variations, $T_f=T_g$ exactly
    when $f=g$ a.e. So it is the \textbf{equivalence classes}, not the
    individual functions, that embed injectively into $\D'$. A distribution of
    the form $T_f$ is called \emph{regular}. The $\delta$ function,
    $\mathrm{p.v.}\,1/x$, and $\delta'$ are not regular.
  \item[Language.] \leavevmode
    \begin{itemize}
      \item ``$T\in\D'$ is in $L^2$'' means $T=T_f$ for some $f\in L^2$.
      \item ``Restricted to $(0,1)$, $T$ is smooth'' refers to its
        \emph{singular support}, the set where $T$ fails to be given by a
        smooth function.
      \item ``The support of $T$'' is the complement of the largest open set
        on which $T=0$.
    \end{itemize}
  \item[Derivatives.] Every distribution has derivatives of every order,
    defined by
    \[
      \pair{T'}{\varphi} = -\pair{T}{\varphi'} .
    \]
    The price is that you \textbf{can't multiply} two distributions in general.
    $\delta\cdot\delta$, for example, is undefined.
\end{description}

\section{The chain of inclusions}

\[
  \Cc \subset C^k \subset C \hookrightarrow L^p_{\mathrm{loc}}
  \subset \Lloc \hookrightarrow \D' .
\]
The arrow $C\hookrightarrow L^p_{\mathrm{loc}}$ is injective on continuous
functions. On arbitrary functions the map to classes is not injective, because
many functions share one class. Sobolev spaces fit in the same way, for example
\[
  H^1_0 \subset L^2 \subset H^{-1} \subset \D',
\]
which is the ``Gelfand triple'' used in weak formulations of PDEs.

\section{Summary table}

\begin{center}
\renewcommand{\arraystretch}{1.3}
\small
\begin{tabular}{>{\raggedright}p{2.6cm} >{\raggedright}p{3.6cm}
                >{\raggedright}p{3.9cm} >{\raggedright\arraybackslash}p{3.9cm}}
  \toprule
  & \textbf{Ordinary functions} & \textbf{Equivalence classes}
  & \textbf{Distributions} \\
  \midrule
  Object & value at every point & functions equal a.e.
    & functional on test functions \\
  Equality & everywhere & almost everywhere & against every $\varphi$ \\
  $f(x_0)$ defined? & yes
    & no (only if a continuous representative exists) & no \\
  Derivative & classical, may not exist & weak, may not exist
    & always exists \\
  Products & yes
    & yes, if the integrability works out (e.g.\ $L^2\cdot L^2\subset L^1$)
    & not in general \\
  Typical spaces & $C$, $C^k$, $C^{k,\alpha}$, $C^\infty$
    & $L^p$, $\Lloc$, $W^{k,p}$, $H^k$
    & $\D'$, $\Sch'$, $\E'$, $H^{-s}$ \\
  Structure & Banach ($C^k$), Fr\'echet ($C^\infty$) & Banach / Hilbert
    & weak-* topology, generally not normed \\
  \bottomrule
\end{tabular}
\end{center}

\section{Three examples where the levels give different answers}

\begin{enumerate}
  \item \textbf{$f(x)=|x|$}
    \begin{itemize}
      \item \emph{Classical:} differentiable everywhere except $x=0$.
      \item \emph{a.e.\ / weak:} $f'=\operatorname{sign}(x)$, an $L^\infty$
        class. The missing value at $0$ doesn't matter.
      \item \emph{Distributional:} $f''=2\delta$, which is not a function.
    \end{itemize}

  \item \textbf{Heaviside function $H$}
    \begin{itemize}
      \item \emph{Classical:} $H'=0$ everywhere except $0$.
      \item \emph{a.e.:} the pointwise a.e.\ derivative is $0$.
      \item \emph{Weak:} $H$ has \textbf{no} weak derivative in $\Lloc$, so
        $H\notin W^{1,1}_{\mathrm{loc}}$.
      \item \emph{Distributional:} $H'=\delta$.
    \end{itemize}
    So ``differentiable a.e.'' is \textbf{not} the same as ``weakly
    differentiable.''

  \item \textbf{Cantor function $c$}
    \begin{itemize}
      \item \emph{Classical / a.e.:} $c$ is continuous, and $c'=0$ a.e.
      \item \emph{Distributional:} $c'$ is the Cantor measure. That measure is
        singular, so it isn't an $\Lloc$ function, and $c\notin W^{1,1}$.
    \end{itemize}
    This is why Sobolev functions in 1D must be \emph{absolutely} continuous:
    continuity alone is not enough.
\end{enumerate}

\end{document}
